\chapter{Introduction}\label{sec:1}

\begin{quote}
\textbf{DRAFT --- not submission-ready.} - Results are final as of
2026-10-12 (all 16 trained runs; Chapter~\ref{sec:4}). - Citation venues
were checked on 2026-10-11; see the reference table in
Chapter~\ref{sec:2}. - Claims about how \emph{other} papers report their
experiments (number of GPUs, number of runs) come from the project notes
and must be checked against each paper before submission. They are
marked ‡.
\end{quote}

\section{Background}\label{sec:1.1}

Large language models (LLMs) can answer a wide range of questions, but
their knowledge is fixed when training ends, and they sometimes state
false information with confidence. A common remedy is
\textbf{retrieval-augmented generation} (RAG): before answering, the
system retrieves relevant documents and conditions its answer on them.
Many questions cannot be answered with a single retrieval. A
\emph{multi-hop} question such as ``Which company formed the boy group
that debuted with the album \emph{2014 S/S}?'' requires finding one fact
(the group), then using it to find another (the company).

This has motivated \textbf{search agents}: LLMs that decide for
themselves when to search, what to search for, and when to stop and
answer (Yao et al., 2023). Prompting alone does not reliably produce
good search behaviour, so recent work trains search agents with
\textbf{reinforcement learning} (RL). In particular, \textbf{Search-R1}
(Jin et al., 2025) trains a model with an outcome reward, meaning the
model is rewarded only when its final answer is correct. It uses
\textbf{GRPO} (Shao et al., 2024), an RL algorithm that compares several
attempts at the same question instead of training a separate value
model. GRPO with rule-checked rewards has also been central to recent
reasoning models (DeepSeek-AI, 2025).

\section{Problem statement}\label{sec:1.2}

Rewarding only the final answer creates a problem. A correct answer does
not show that the agent found, used, or could point to the evidence
behind it. In RL terms, the answer reward is a \textbf{proxy} for what
is actually wanted: a correct answer \emph{supported by evidence}.
Optimizing a proxy invites \textbf{reward hacking}, where the policy
finds behaviour that scores well on the proxy but not on the real goal
(Skalse et al., 2022). This thesis addresses four related gaps:

\begin{itemize}
\tightlist
\item
  \textbf{G1, correct is not grounded.} An agent trained with
  answer-only reward can learn to answer without showing evidence, or to
  search in wasteful ways, as long as the final answer is right. Recent
  work reports such shortcut behaviour in outcome-trained search agents
  (Zhang et al., 2026).
\item
  \textbf{G2, evidence rewards are expensive.} Existing rewards for
  citations and evidence typically use an LLM as a judge (Zhang et al.,
  2026; Xu et al., 2026). That adds compute cost and latency, and the
  judge can be biased or wrong.
\item
  \textbf{G3, a compute barrier.} RL for search agents is usually
  reported at multi-GPU scale‡. It is unclear whether the same effects
  can be studied cheaply.
\item
  \textbf{G4, weak evidence about what matters.} RL results vary
  considerably between random seeds (Henderson et al., 2018). Papers
  that report single runs‡ cannot separate the effect of a method from
  seed noise, nor separate the effect of the reward from the effect of
  the optimizer (Agarwal et al., 2021).
\end{itemize}

HotpotQA (Yang et al., 2018) offers a way to address G1 and G2 together.
Its questions are annotated with the exact sentences that support the
answer. If an agent cites sentences, the citations can be checked by a
simple rule against these gold labels, with no judge.

\section{Research questions}\label{sec:1.3}

\begin{itemize}
\tightlist
\item
  \textbf{RQ1.} Does GRPO training improve a 7B search agent on HotpotQA
  over the un-tuned model?
\item
  \textbf{RQ2.} Does a verifiable evidence-grounded reward (VEGR)
  prevent the shortcut behaviour of answer-only reward, without losing
  accuracy?
\item
  \textbf{RQ3.} How do reward variants (gated penalties, a larger
  coverage weight) and optimizer fixes (Dr.~GRPO advantages, dynamic
  sampling) affect grounding and search behaviour?
\end{itemize}

\section{Objectives}\label{sec:1.4}

\begin{enumerate}
\def\labelenumi{\arabic{enumi}.}
\tightlist
\item
  Build a search environment and agent in which every citation can be
  verified exactly against HotpotQA's supporting facts.
\item
  Design a judge-free reward, VEGR, that rewards correct answers,
  correct citations and useful searches, and penalizes invalid actions.
\item
  Train the agent with GRPO and LoRA on a single GPU, using both
  answer-only reward and VEGR, with three random seeds each.
\item
  Compare the trained agents on a frozen holdout set in terms of both
  accuracy and grounding, using paired bootstrap confidence intervals.
\item
  Test whether reward variants and optimizer fixes change grounding and
  stability.
\end{enumerate}

\section{Contributions}\label{sec:1.5}

\begin{enumerate}
\def\labelenumi{\arabic{enumi}.}
\tightlist
\item
  \textbf{VEGR}, a judge-free, sentence-level reward for training search
  agents. Every term is checked by rule against gold supporting facts.
  VEGR also enforces the action protocol: invented citations, malformed
  actions, and running out of budget all receive a −1 penalty. The
  individual reward terms are not new; the contribution is the
  combination and the controlled evaluation of it.
\item
  \textbf{Evidence that answer-only GRPO produces shortcuts.} On holdout
  evaluation, two of three outcome-only seeds stopped citing entirely
  (citation precision 0.00), and the third repeated searches in 229 of
  its 278 completed episodes. The VEGR runs reached the same accuracy
  (+1.0 EM points, 95\% CI {[}−1.9, +3.8{]}) but \textbf{24.4 points
  higher grounded EM} (CI {[}+20.9, +28.1{]}), with citation precision
  at 0.73--0.75 on every seed.
\item
  \textbf{A controlled multi-seed comparison} with paired bootstrap
  confidence intervals, which separates reward design from optimizer
  fixes. It shows that:

  \begin{itemize}
  \tightlist
  \item
    a single-seed improvement of +6 EM points (gated penalties, R2) did
    not replicate across seeds;
  \item
    doubling the coverage weight (R3) raises evidence coverage but
    lowers grounded EM by 8.0 points;
  \item
    Dr.~GRPO advantages (R4) make no measurable difference;
  \item
    dynamic sampling (R0-dyn, one seed) raises outcome-only accuracy by
    8 points but does not prevent the citation shortcut. The shortcut
    comes from the reward, not only from the optimizer.
  \end{itemize}
\item
  \textbf{An open, reproducible single-GPU pipeline.} One training run
  plus evaluation costs about \$2.25 on one 48 GB GPU, and splits,
  configs, logs and per-episode results are all released in the
  repository.
\end{enumerate}

\section{Scope and limitations}\label{sec:1.6}

This thesis is deliberately narrow:

\begin{itemize}
\tightlist
\item
  \textbf{Data and retrieval.} It uses HotpotQA in the distractor
  setting. The agent searches only the 10 paragraphs that come with each
  question, using BM25, so this is not open-web or open-corpus search.
\item
  \textbf{Model.} The model is Qwen2.5-7B-Instruct adapted with LoRA.
\item
  \textbf{Scale.} Each training run uses 256 questions for one epoch,
  and evaluation uses a frozen 300-question holdout. The whole project
  used a compute budget of about \$45.
\item
  \textbf{What is not claimed.} The thesis does not propose a new RL
  algorithm, does not claim state-of-the-art accuracy, and does not
  claim that the findings transfer to open-web search or larger models.
\item
  \textbf{Known loophole.} The environment accepts an answer with no
  citations. This was kept on purpose so that all runs are comparable,
  and its exploitation is reported as a finding.
\end{itemize}

Chapter~\ref{sec:5} discusses the further limitations.

\section{Thesis structure}\label{sec:1.7}

\begin{itemize}
\tightlist
\item
  \textbf{Chapter~\ref{sec:2}} reviews the background: LLMs, retrieval,
  multi-hop QA and citation evaluation, LLM agents, RL for LLMs (PPO,
  GRPO and its fixes, LoRA), RL-trained search agents, evidence rewards,
  and reward hacking and RL evaluation. It ends with the research gaps.
\item
  \textbf{Chapter~\ref{sec:3}} describes the methodology: data splits,
  the search environment, the VEGR reward and its variants, the policy
  model, the GRPO training algorithm, the evaluation protocol, and the
  experiment matrix.
\item
  \textbf{Chapter~\ref{sec:4}} reports the results by research question,
  diagnoses the training dynamics, presents example traces, and
  discusses threats to validity.
\item
  \textbf{Chapter~\ref{sec:5}} concludes, answers the research
  questions, and outlines future work.
\end{itemize}
